
\subsection{Verrouillage à deux phases}


\begin{frame}{Verrouillage à deux phases}

Le principal (et le plus ancien) algorithme pour assurer la
sérialisabilité.

\vfill

Est réputé engendrer beaucoup de blocages et de rejets.

\vfill

Repose sur le verrouillage (surprise!) des nuplets.

\end{frame}


\subsection{Verrouillage}

\begin{frame}{Les verrous}

Pose de \alert{verrous} sur les nuplets.


\begin{redbullet}

   \item Le verrou \alert{partagé} (\textit{shared lock})
          autorise la pose d'autres verrous partagés sur le même tuple. 

   \item Le verrou \alert{exclusif} (\textit{exclusive lock})
           interdit la pose de tout autre verrou, 
           exclusif ou partagé, et donc de toute lecture ou écriture par une autre transaction.
\end{redbullet}

\vfill

Verrous = \alert{blocages}. Il faut
en poser

\begin{blocgauche}
\begin{redbullet}
   \item  \alert{le moins possible}, surtout pour les verrous exclusifs.
  \item \alert{pour la durée la plus courte possible}
\end{redbullet}
\end{blocgauche}

\end{frame}

\begin{frame}{Transactions et verrous}

Le contrôleur pose des verrous pour une transaction $T$.

\begin{redbullet}
  \item On ne peut poser un \alert{verrou partagé} sur un tuple $a$ 
       que s'il n'y a que des verrous partagés sur ce tuple.
  \item On ne peut poser un \alert{verrou exclusif} que si
  
     \begin{itemize}
        \item il n'y a aucun autre verrou, \alert{ou},
         \item il y a un verrou partagé déjà posé par $T$ elle-même (\textit{upgrade})
         \end{itemize}
 \end{redbullet}
 
 \vfill
\begin{blocgauche}
\begin{redbullet}
  \item Une transaction qui n'obtient pas un verrou,  \alert{est mise en attente}. 
\item Les verrous ne sont libérés qu'au moment du \texttt{commit} ou
\texttt{rollback}.
\end{redbullet}
\end{blocgauche}

\end{frame}


\begin{frame}{L'algorithme}

On pose des verrous pour empêcher l'apparition de cycles.


\begin{redbullet}
	\item \alert{Lecture} sur $x$,
            on regarde s'il y a un verrou \alert{exclusif} sur $x$\,;
	\begin{itemize}
		\item si oui la transaction $T_i$ 
	                est mise en attente.
		\item si non, la transaction pose un verrou en lecture
		       et l'opération est exécutée.
	\end{itemize}
	\item \alert{Ecriture} sur $x$,
              on regarde s'il y a un verrou \alert{quelconque}  sur $x$\,;
	\begin{itemize}
		\item si oui la transaction $T_i$ 
	                est mise en attente.
		\item si non, la transaction pose un verrou en écriture
		       et l'opération est exécutée.
	\end{itemize}
\end{redbullet}

\vfill

\begin{blocgauche}
\alert{Important}\,: ne fonctionne que si les verrous ne sont pas
relâchés avant le \texttt{commit} ou le \texttt{rollback}.
\end{blocgauche}
\end{frame}


\begin{frame}{Exemple}

Prenons l'exécution concurrente\,:  $r_1[x] w_2[x] w_2[y] C_2 w_1[y] C_1$

\vfill

\begin{redbullet}
  \item $T_1$ pose un verrou partagé sur $x$, lit $x$
          mais ne relâche pas le verrou\,;
  \item $T_2$ tente de poser un verrou exclusif sur $x$\,:
              impossible puisque $T_1$ détient un verrou partagé,
                 {\em donc $T_2$ est mise en attente}\,;
  \item $T_1$ pose un verrou exclusif sur $y$, modifie $y$,
            valide\,;  ses verrous  
            sont relâchés\,;
  \item $T_2$ est libérée\,: elle pose un verrou exclusif sur $x$, et 
              le modifie\,;
  \item $T_2$ pose un verrou exclusif sur $y$, et modifie $y$\,;
  \item $T_2$ valide, ce qui relâche les verrous sur $x$ et $y$.
\end{redbullet}

Exécution après réordonnancement: $r_1[x] w_1[y] w_2[x] w_2[y]$
\end{frame}


%%%%%%%%%%%%%%%%%%%%%%%%%%%%%%%%%%%%%%%%%%%%%%%%%%%
\begin{frame}{Un gros problème\,: les \texttt{deadlock}}


Reprenons notre exemple des mises à jour perdues.

\begin{center}
$r_1(s)  r_1(c_1) r_2(s)  r_2(c_2) 
 w_2(s)  w_2(c_2) C_2 w_1(s) w_1(c_1) C_1$
\end{center}

\begin{redbullet}
  \item $T_1$ lit $s$ et $c_1$, qui sont verrouillés en lecture.
  \item $T_2$ lit $s$ et $c_2$, qui sont verrouillés en lecture\,;
     \alert{$s$ partage deux verrous en lecture}

  \item $T_2$ veut écrire $s$\,: conflit, donc blocage de $T_2$.

  \item $T_1$ veut écrire $s$\,: conflit, donc blocage de $T_1$.
\end{redbullet}

C'est \alert{l'étreinte fatale} (\textit{deadlock})\,!!

\vfill

\begin{blocgauche}
 Le système va rejeter une
des transactions.
Très regrettable, mais encore mieux que d'introduire des anomalies (?)
\end{blocgauche}
\end{frame}


%%%%%%%%%%%
\begin{frame}{À retenir}

Le 2PL est le seul algorithme actuellement utilisé pour garantir la sérialisabilité.

\vfill

Il repose intensivement sur la pose de verrous qui ne sont relâchés qu'à la fin de la transaction.

\vfill
\begin{blocgauche}
Il entraîne des interblocages, et donc le rejet de certaines transactions\,: difficilement explicable pour
un utilisateur, pose le problème de resoumettre la transaction.
\end{blocgauche}

\end{frame}
